% basic nastaveni
\documentclass[12pt]{article}
\setlength{\parindent}{0pt}
\setlength{\parskip}{1em}
\usepackage[utf8]{inputenc}
\usepackage[czech]{babel}
\usepackage[T1]{fontenc}
\usepackage{graphics}
\usepackage{caption}
\usepackage{hyperref}
\usepackage{dsfont}
\usepackage{color}
\usepackage{float}
\usepackage{enumitem}
\usepackage{graphicx}

% matematicke package
\usepackage{tikz-cd}
\usetikzlibrary{babel}
\usepackage{amsmath}
\usepackage{amssymb}
\usepackage{amsfonts}
\usepackage{amssymb}
\usepackage{geometry}
\geometry{margin=2.5cm}
\usepackage{pgfplots}
\pgfplotsset{compat=1.18}
\usepackage{mathrsfs}
\usetikzlibrary{shapes.misc}
\usetikzlibrary{angles,quotes}

\title{Maturitní otázka číslo 9: Kombinatorika}
\author{}
\date{}

\begin{document}
\maketitle

\section*{Základní kombinatorická pravidla}
\subsection*{Kombinatorické pravidlo součinu}
Definice: Počet všech uspořádaných k-tic, jejichž první člen lze vybrat $n_1$ způsoby, druhý člen po výběru prvního členu $n_2$ způsobu atd. až k-tý člen po výběru všech předcházejících členů $n_k$ způsoby, je roven $n_1 \cdot n_2 \cdot ... \cdot n_k$.

Příklad: U stánku nabízejí tři různé zmrzliny a čtyři různé polevy. Kolik druhů zmrzliny s polevou můžeme vytvořit, pokud nebudeme míchat polevy ani druhy zmrzliny?

Řešení: Ke každému ze tří druhů zmrzliny můžeme přiřadit čtyři polevy, můžeme tedy vytvořit $3 \cdot 4 = 12$ různých zmrzlin s polevou.

Převedeno na pojmy z definice bude $k = 2, n_1 = 3, n_2 = 4, n_1\cdot n_2 = 12$.

\subsection*{Kombinatorické pravidlo součtu}
Definice: Jsou-li $A_1, A_2, ..., A_n$ konečné množiny, které mají po řadě $p_1, p_2, ..., p_n$ prvků, a jsou-li každé dvě disjunktní, pak počet prvků množiny $A_1 \cup A_2 \cup ... \cup A_n$ roven $p_1 + p_2 + ... + p_n$.

\section*{Faktoriál}
Faktoriál se značí $!$ za číslem. Faktoriál $0! = 1$, pro středoškolskou matematiku je faktoriál definovaný pouze na přirozených číslech s $0$
\section*{Variace bez opakování}
Definice: k-členná variace z $n$ prvků je uspořádaná k-tice sestavená z těchto prvků tak, že každý se v ní vyskytuje nejvýše jednou.
\[V(k,n) = \frac{n!}{(n-k)!}\]
Příklad: Závod běžel Petr, Roman, Simona a Jamal. Kolika různými způsoby mohli obsadit stupně vítězů?

Řešení: Na prvním místě mohl doběhnout kterýkoliv z těch čtyř, tudíž máme 4 možnosti.
 Na druhé místo můžeme dosadit už jen 3 různé, protože jeden už je na prvním místě, na třetí místo máme na výběr jen ze 2\\
Podle pravidla kombinatorického součinu je tedy $4 \cdot 3 \cdot 2 = 24$ možností, jak mohli doběhnout. V tomto příkladě bude $k = 3$ a $n = 4$

\section*{Permutace}
Permutace je zvláštní případ variace, kde $k=n$, tudíž ze zadaných prvků postupně vybereme všechny.

Definice: Permutace z $n$ prvků je uspořádaná $n-tice$ sestavená z těchto prvků tak, že každý se v ní vyskytuje právě jednou\\
\[P(n) = n!\]

Příklad: Určete kolika způsoby se může $9$ žáků seřadit tak aby Milan a Lumír nestáli
 vedle sebe.

Řešení: Počet celkových seřazení je $9!$ (i těch kde stojí Lumír a Milan vedle sebe). 
Počet způsobů jak je seřadit tak, aby Lumír a Milan stáli vedle sebe je $8!$, protože 
je prostě spojíme do dvojice (jednoho prvku). V každém seřazení se ale ještě ti dva 
mohou prohodit a pořád budou stát vedle sebe, tudíž celkový počet možností, kde stojí 
vedle sebe je $2\cdot 8!$. Výsledek tedy bude: $9! - 2\cdot8! = 9 \cdot 8! - 2\cdot 8! = 7\cdot 8!$

\section*{Kombinace}
U kombinací nezáleží na pořadí prvků, pokud tedy vybíráme z množiny $A = \{1,2,3,4\}$ tři prvky, kombinace $\{1,2,3\}$ a $\{3,2,1\}$ jsou stejné (jsou vybrané stejné prvky)

Definice: $k$-členná kombinace z $n$ prvků je neuspořádaná $k$-tice sestavená z těchto prvků tak, že každý se v ní vyskytuje nejvýše jednou.
\[K(k,n) = \frac{n!}{k!(n-k)!}\]
Příklad: Ve třídě je 26 žíáků, kolika způsoby můžeme vybrat 2 zástupce třídy

Řešení: $K(2,26) = \frac{26!}{2!\cdot24!} = 325$

\subsection*{Kombinační číslo}
Kombinační číslo se zapisuje jako $\binom{n}{k}$, čteme ho jako $n$ nad $k$.

Platí také tento vztah pro všechna nezáporná čísla $n, k, k\leq n\quad$ $\binom{n}{n-k} = \binom{n}{k}$ 
To proto, že pokud vybíráme například z $10$ prvkové množiny $4$ prvky, tak těch 
zbývajících $6$ bude stejných, jako kdybychom vybírali $6$ prvků.

Důkaz: \[\binom{n}{n-k} = \frac{n!}{(n-k)![n-(n-k)]!} = \frac{n!}{(n-k)!k!} = \binom{n}{k}\]

Vlastnosti kombinačních čísel:
\[\binom{n}{0} = \binom{n}{n} = 1\]
\[\binom{n}{1} = n\]
\[\binom{n}{k}+ \binom{n}{k+1} = \binom{n+1}{k+1} \quad \text{Pascalovo pravidlo}\]

\textbf{Důkaz Pascalova pravidla}
\[\binom{n}{k} + \binom{n}{k+1} = \frac{n!}{(n-k)!k!} + \frac{n!}{(n-(k+1))!(k+1)!} = \frac{n!(k+1)!}{(k+1)!(n-k)!} + \frac{n!(n-k)}{(k+1)!(n-k)!}=\]
\[=\frac{n!(k+1)! + n!(n-k)}{(k+1)!(n-k)!} = \frac{n!(n+1+k-k)}{(k+1)!(n-k)!} = \frac{(n+1)!}{(k+1)!(n+1-(k+1))!} = \binom{n+1}{k+1} \]

\section*{Variace s opakováním}
Prvky se mohou ve výběru opakovat a záleží na pořadí

Definice: $k$-členná variace s opakováním z $n$ prvků je uspořádaná $k$-tice sestavená z těchto prvků tak, že každý se v ní vyskytuje nejvýše $k$-krát.
Vzorec: $V'(k,n) = n^k$

Příklad: Vyberme 4 pastelky z krabice se žlutými, z krabice s modrými a z krabice s červenými. Kolik různých čtveřic pastelek můžeme vytvrořit, pokud je dáme za sebe na stůl? Na první místo můžeme dát 3 různé pastelky, na druhé místo také, protože se mohou opakovat atd.
Můžeme tedy vytvořit podle kombinatorického pravidla součinu $3^4 = 81$ různých pořadí pastelek, které se mohou opakovat.

\section*{Permutace s opakováním}
Definice: Permutace s opakováním z $n$ prvků je uspořádaná $k$-tice sestavená z těchto prvků tak, že každý se v ní vyskytuje alespoň jednou.

Přirozené číslo $n$ udává počet různých prvků.\\
Jednotlivé prvky se mohou opakovat a je zvykem označovat počet opakování prvního prvku $k_1$, počet opakování druhého $k_2$ atd. Přirozené číslo $k$ označuje počet všech prvků, jejichž různá pořadí zkoumáme, proto platí $k = k_1 + k_2 + \dots + k_n$.

Příklad: Mějme 3 krabice s pastelkami(žlutá, modrá, červená) a permutaci, ve které se žlutá bude opakovat $k_1 = 4$, modrá $k_2 = 3$ a červená $k_3 = 5$. Počet různých variant, pokud by každá pastelka byla rozdílná by byl $(k_1 + k_2 + k_3)! = 12!$ (variace). Počet $k_1$ pastelek které lze na 4 místa na kterých jsou doplnit různými způsoby je celkem $k_1! = 4!$, modré jde doplnit na jejich místo celkem $k_2 = 3!$ způsoby atd.\\
Počet různých permutací s opakováním bude tedy $\frac{(4 + 3 + 5)!}{4!3!5!}$

Obecný vzorec je tedy:
\[P'(k_1, k_2, \dots, k_n) = \frac{(k_1 + k_2 +\dots + k_n)!}{k_1!k_2!\dots k_n!}\]

\section*{Kombinace s opakováním}
Definice: $k$-členná kombinace s opakováním z $n$ prvků je neuspořádaná $k$-tice sestavená z těchto prvků tak, že každý se v ní vyskytuje nejvýše $k$-krát.

Příklad: Určete kolika způsoby je možné rozmístit 3 kuličky do čtyř krabiček.

Řešení: Třikrát vybíráme jednu ze čtyř krabiček, jde tedy o tříčlennou kombinaci s opakováním ze čtyř prvků. $k = 3, n = 4$
Jedna z možných kombinací tedy bude $.|..||$, kde $|$ odděluje jednotlivé krabičky od sebe (1. krabička má 1 kuličku, 2. krabička má 2 a třetí je prázdná). Počet možných pořadí uvedených 6 symbolů můžeme určit pomocí permutace $P'(3,3) = \frac{(3+3)!}{3!3!} = 20$\\
Pokud zobecníme úlohu tak, že budeme hledat počet způsobů jak rozmísit $k$ stejných kuliček do $n$ krabiček, byly by to $(k+n-1)$-tice ve kterých je tečka $k$-krát a symbol $|$ $(n-1)$-krát (počet oddělení krabiček). Můžeme to zapsat jako permutaci 
\[P'(k,n-1) = \frac{(k+n-1)!}{k!(n-1)!} = \frac{n+k-1}{k!(n-1)!} = \binom{n+k-1}{k}\]\\
To je výsledný vzorec pro kombinace s opakováním, které se značí $K'(k,n)$

\section*{Binomická věta}

\end{document}